\documentclass[12pt]{article}
\usepackage[utf8]{inputenc}
\usepackage[T1]{fontenc}
\usepackage{hyperref}
\title{Tutsa Naga Language Resources for K-12 Curriculum Development}
\author{Compiled by Claude Sonnet 5 (claude-sonnet-5) for Andr\'as Kornai}
\date{}
\begin{document}
\maketitle

\section{Introduction}
Tutsa Naga (ISO 639-3 tvt, Glottocode tuts1235) is a Sino-Tibetan language. Glottolog places it in the Tutsic branch of Noctean, within North Patkaian. Search results place its speakers in southern Changlang and eastern Tirap districts of Arunachal Pradesh and in Tinsukia district of Assam. About 25,000 speakers are cited (2001). The name ``Tangsa'' is used as an umbrella term for related varieties, so sources need care.

\section{Primary Educational Resources}
None were found in search results.

\section{Linguistic and Reference Materials}
\begin{itemize}
\item Wikipedia, Tutsa language (a stub): \url{https://en.wikipedia.org/wiki/Tutsa_language}
\item Glottolog entry: \url{https://glottolog.org/resource/languoid/id/tuts1235.bigmap.html}
\end{itemize}
No dictionary or grammar turned up in the searches.

\section{Religious and Community Resources}
\begin{itemize}
\item Scripture Earth listing for Tutsa (links to scripture on YouVersion and eBible.org): \url{https://www.scriptureearth.org/00eng.php?iso=tvt}
\end{itemize}

\section{Digital and Audio Resources}
\begin{itemize}
\item Global Recordings Network, Tutsa recordings (including a Words of Life program): \url{https://globalrecordings.net/en/language/tvt}
\end{itemize}

\section{Curriculum Development Context}
The language is not known to be taught in schools. The search found only scripture, audio evangelism material and catalogue entries. A curriculum would need new primers and a community-agreed orthography.
\end{document}